% 原书第57--62页；2.3全部五分目.

\section{\texorpdfstring{$X_0(N)$}{X₀(N)} 的典范模型}\label{sec:modular-2-3}

我们知道，每个紧Riemann曲面都可由代数方程定义.下面将这一说法稍加形式化.



\subsection{代数簇}\label{sec:modular-2-3-1}

设 $k$ 为域.$k$ 上的\newterm{modular-algebraic-variety}{代数簇}{algebraic variety} $X$ 是一个有限型整 $k$ 概形.因而它是一个局部同构于有限生成整 $k$ 代数的谱的赋环空间.簇的态射就是赋环空间的态射.



若一个簇同构于 $\mathbb P^n$ 的闭子簇（由齐次多项式定义），则称它为射影簇.簇有唯一的\newterm{modular-generic-point}{普遍点}{generic point}，该点的局部环是 $X$ 的有理函数域 $k(X)$.$X$ 的维数（即不可约闭集的极大链的长度）等于 $k(X)$ 在 $k$ 上的超越次数.



$k$ 上的曲线 $C$ 是维数为 $1$ 的簇.若 $k$ 为完美域，且 $C$ 的局部环都是离散赋值环，则称 $C$ 在 $k$ 上光滑.可以证明，每条``紧合''（或``完备''）曲线实际上都是射影曲线.代数曲线理论的最初结果之一是，对应
\begin{equation*}
\begin{tikzcd}
C \arrow[r,mapsto] & k(C)
\end{tikzcd}
\end{equation*}
在 $k$ 上代数曲线（以非常值态射为态射）与 $k$ 上有限生成、超越次数为 $1$ 的域之间给出范畴的反等价.\footnote{译注：此处按原文保留.这一反等价通常须将曲线限制为光滑完备曲线；下文给出的重构正是完备模型的重构.}下面说明如何从这样的域 $K$ 重构曲线 $C_K$.取其底空间为
\begin{align*}
|C_K|=\{\eta\}\cup\{\text{离散赋值环 }k\subset V\subset K\mid\operatorname{Frac}(V)=K\},
\end{align*}
其中 $\eta$ 将作为普遍点.若一个集合是 $C_K$ 中某个有限子集的补集，并且含有 $\eta$，就规定它为开集.对于开集 $U$，置 $\mathcal O_{C_K}(U)=\bigcap_{V\in U}V\subset K$.可以证明，这是有限生成的 $k$ 代数，也是Dedekind环.\footnote{译注：原文对任意开集如此断言.对全空间 $U=C_K$，所述有限生成性与Dedekind环断言须区别处理；此处未改写原文的重构过程.}预层 $U\mapsto\mathcal O_{C_K}(U)$ 是层，因而在 $C_K$ 上定义赋环空间结构；并可验证，$U$ 上的诱导结构就是 $\operatorname{Spec}(\mathcal O_{C_K}(U))$ 的结构.于是 $C_K$ 是一条曲线，并且容易验证它是完备的.



\subsection{一维复环面的代数化}\label{sec:modular-2-3-2}

紧Riemann曲面的一种实际的代数化方法，是取其函数域所对应的光滑完备代数曲线.但对于环面 $\mathbb C/\Lambda$，可以显式写出全纯依赖于 $\Lambda$ 的方程.考虑复变量 $x\in\mathbb C$ 的级数
\begin{align*}
\wp'_\Lambda(x)=-2\sum_{\omega\in\Lambda}\frac1{(x-\omega)^3}.
\end{align*}
在 $\Lambda$ 的任意基本域上，去掉在其中有极点的有限个项以后，此级数正规收敛.因此，该函数是 $\Lambda$ 周期的，在 $\Lambda$ 之外全纯，并在每个 $\omega\in\Lambda$ 处有三阶极点.它有原函数，即变量 $x\in\mathbb C$ 的Weierstrass函数 $\wp_\Lambda$\footnote{译注：原书的求和指标写作 $\omega\in\Lambda$，未排除 $\omega=0$；通常此处应对 $\Lambda\setminus\{0\}$ 求和.}：
\begin{align*}
\wp_\Lambda(x)=\frac1{x^2}+\sum_{\omega\in\Lambda}\left(\frac1{(x-\omega)^2}-\frac1{\omega^2}\right).
\end{align*}
它同样是 $\Lambda$ 周期的亚纯函数，但在 $\omega\in\Lambda$ 处有二阶极点.因此，$\wp_\Lambda$ 与 $\wp'_\Lambda$ 下降为 $\mathbb C/\Lambda$ 上的亚纯函数.考察 $\mathbb C/\Lambda$ 上一般亚纯函数可能具有的极点，可以证明 $\mathcal M_{\mathbb C/\Lambda}=\mathbb C(\wp_\Lambda,\wp'_\Lambda)$.用类似思路，可证明函数
\begin{align*}
\wp_\Gamma^{\prime 2}-4\wp_\Lambda^3+g_2(\Lambda)\wp_\Lambda+g_3(\Lambda)
\end{align*}
是全纯的，因而是常数，进而恒为零.\footnote{译注：原书第一项的下标确为 $\Gamma$，其余项均为 $\Lambda$；此处保留原式.}这给出域 $\mathcal M_{\mathbb C/\Lambda}$ 的一个表示，并表明映射
\begin{equation*}
\begin{tikzcd}[column sep=large,row sep=0.2em]
\mathbb C/\Lambda \arrow[r] & \mathbb P^2(\mathbb C)\\
x\ne0 \arrow[r,mapsto] & {[\wp(x):\wp'(x):1]}\\
0 \arrow[r,mapsto] & {[0:1:0]}
\end{tikzcd}
\end{equation*}
是从 $\mathbb C/\Lambda$ 到椭圆曲线 $y^2=4x^3-g_2x-g_3$ 所对应Riemann曲面的同构.注意，让 $\Lambda$ 按 $\Lambda=\Lambda_z=z\mathbb Z\oplus\mathbb Z$（$z\in\mathbb H$）变化，便得到 $\mathbb H$ 上一族解析的椭圆曲线，它位于 $\mathbb H\times\mathbb P^2(\mathbb C)$ 内，非齐次方程为 $y^2=4x^3-g_2(z)-g_3(z)$.\footnote{译注：原书此式 $g_2(z)$ 后没有因子 $x$，此处保留.}这一族的判别式是模形式 $\Delta\in\mathcal S_{12}(\Gamma(1))$，而其 $j$ 不变量就是模函数 $j$.



回忆，$Y(1)$ 在同构意义下参数化复环面.自然会问：上述族是否``下降''为 $Y(1)$ 上的椭圆曲线族？甚至还希望这一族代数地依赖于 $Y(1)$ 上的坐标 $j$，且曲线 $E_j$ 的不变量就是 $j$……但由于椭圆点对应于有``过多自同构''的曲线，这样的要求有些过高.



因此，把上述族限制在普通点轨迹上.注意，这正是 $g_2\ne0$ 且 $g_3\ne0$ 的轨迹，也就是 $j\ne0,1728$ 的轨迹.作变量替换
\begin{align*}
(x,y)=\bigl((g_3/g_2)x',\,(g_3/g_2)^{3/2}y'\bigr),
\end{align*}
可将限制后的族写成
\begin{align*}
E_z:\quad y^2=4x^3-\frac{27j(z)}{j(z)-1728}x-\frac{27j(z)}{j(z)-1728}.
\end{align*}
这显然下降为以 $j$ 为坐标、删去 $j=0,1728$ 两点的仿射直线上的代数椭圆曲线族，其方程为
\begin{align*}
E_j:\quad y^2=4x^3-\frac{27j}{j-1728}x-\frac{27j}{j-1728}.
\end{align*}
我们知道，椭圆曲线可在 $\mathbb Q$ 上定义，当且仅当其 $j$ 不变量为有理数；这便说明，取函数域为 $\mathbb Q(j)$ 的射影直线 $\mathbb P^1$ 作为 $X(1)$ 在 $\mathbb Q$ 上的模型 $X(1)_{\mathbb Q}$ 是合理的.



\subsection{\texorpdfstring{$X_0(N)$}{X₀(N)} 的代数化及有理模型}\label{sec:modular-2-3-3}

下面描述 $X_0(N)$ 的亚纯函数域.我们知道，$j$ 给出亚纯函数 $j\in\mathcal M_{X_0(N)}$.由前面见过的计算，亚纯函数 $j_N(z)=j(Nz)=j(\begin{psmallmatrix}N&0\\0&1\end{psmallmatrix}z)$ 也是级为 $\Gamma_0(N)$ 的模函数.

\begin{theorem*}

有 $\mathcal M_{X_0(N)}=\mathbb C(j,j_N)$.此外，存在对称多项式 $f_N\in\mathbb Z[X,Y]$，使得 $f_N(X,Y)=0$.\footnote{译注：原文确实写作 $f_N(X,Y)=0$；按下文应为 $f_N(j,j_N)=0$.}最后，$f_p(X,Y)\equiv(Y-X^p)(Y^p-X)\pmod p$.

\end{theorem*}

\begin{proof}

回忆，我们已经知道 $\mathcal M_{X_0(N)}$ 在 $\mathcal M_{X(1)}=\mathbb Q(j)$ 上的次数为 $[\Gamma(1):\Gamma_0(N)]$.\footnote{译注：原书此处写 $\mathbb Q(j)$；复Riemann曲面 $X(1)$ 的亚纯函数域应为 $\mathbb C(j)$，下文也采用后者.}



选取 $\Gamma_0(N)\backslash\Gamma(1)$ 的代表元集 $\mathcal L$.每个 $\alpha\in\Gamma(1)$ 诱导置换 $\gamma\mapsto\gamma_\alpha$，由下式确定：
\begin{align*}
\Gamma_0(N)\gamma\alpha=\Gamma_0(N)\gamma_\alpha.
\end{align*}
对任意函数 $f\in\mathcal M_{\Gamma_0(N)}$，考虑多项式
\begin{align*}
P_f(Y):=\prod_{\gamma\in\mathcal L}(Y-f(\gamma z))=\sum_{i=0}^{|\mathcal L|}a_i(z)Y^i.
\end{align*}
系数 $a_i(z)$ 是各 $f(\gamma z)$（$\gamma\in\mathcal L$）的对称多项式.因此，对每个 $\alpha\in\Gamma(1)$ 都有 $a_i(\alpha z)=a_i(z)$，即 $a_i(z)\in\mathcal M_{\Gamma(1)}=\mathbb C(j)$，从而 $P_f\in\mathbb C(j)[Y]$.



另一方面，考虑 $f$ 在 $\mathbb C(j)$ 上的极小多项式 $Q_f(Y)=\sum_{j=0}^r b_j(z)Y^j$（故 $b_j\in\mathbb C(j)$）.由于 $b_j$ 对每个 $\alpha\in\Gamma(1)$ 都不变，函数 $f(\gamma z)$（$\gamma\in\mathcal L$）也是此多项式的根.于是 $Q_f$ 在 $\mathbb C(j)[Y]$ 中整除 $P_f$.



将此应用于函数 $j_N$.首先注意，函数 $j_N(\gamma z)$（$\gamma\in\mathcal L$）两两不同.事实上，若 $j(N\gamma z)=j(N\gamma'z)$，则知道存在 $\alpha\in\Gamma(1)$，使得对每个 $z$ 都有 $N\gamma z=\alpha N\gamma'z$.取在 $\Gamma(1)$ 中稳定子为 $\{\pm1\}$ 的 $z$（即普通点），便得
\begin{align*}
\begin{pmatrix}N&0\\0&1\end{pmatrix}\gamma
&=\pm\alpha\begin{pmatrix}N&0\\0&1\end{pmatrix}\gamma',\\
\gamma\gamma'^{-1}
&=\begin{pmatrix}N&0\\0&1\end{pmatrix}^{-1}(\pm\alpha)\begin{pmatrix}N&0\\0&1\end{pmatrix}.
\end{align*}
特别地，
\begin{align*}
\gamma\gamma'^{-1}\in\Gamma(1)\cap\begin{pmatrix}N&0\\0&1\end{pmatrix}^{-1}\Gamma(1)\begin{pmatrix}N&0\\0&1\end{pmatrix}=\Gamma_0(N),
\end{align*}
所以 $\gamma=\gamma'$.因此，$j_N$ 在 $\mathbb C(j)$ 上的极小多项式 $Q_{j_N}$ 等于
\begin{align*}
P_{j_N}(Y)=\prod_{\gamma\in\mathcal L}(Y-j_N(\gamma z))\in\mathbb C(j)[Y].
\end{align*}
其次数等于 $[\mathcal M_{X_0(N)}:\mathbb C(j)]$，于是 $\mathcal M_{X_0(N)}=\mathbb C(j,j_N)$.



现在，$P_{j_N}(Y)$ 的系数既是 $j$ 的有理函数，又是 $z$ 的全纯函数.由于 $j:\mathbb H\longrightarrow\mathbb C$ 是满射，这些系数必为 $j$ 的多项式（否则作为 $z$ 的函数会有极点）.因此，对唯一的多项式 $f_N(X,Y)\in\mathbb C[X,Y]$，有 $P_{j_N}(Y)=f_N(j,Y)$.



下面证明 $f_N(X,Y)$ 关于 $X,Y$ 对称.将 $z$ 换为 $-1/(Nz)$，可得 $f_N(j(-1/(Nz)),j(-1/z))=0$.由 $j$ 对 $\begin{psmallmatrix}0&-1\\1&0\end{psmallmatrix}$ 的不变性，$j(-1/(Nz))=j(Nz)$ 且 $j(-1/z)=j(z)$，故 $f_N(j_N,j)=0$.由于 $f_N(X,Y)$ 在 $\mathbb C(X)[Y]$ 中不可约，$f_N(Y,X)$ 在 $\mathbb C(X)[Y]$ 中是 $f_N(X,Y)$ 的倍式，实际上在 $\mathbb C[X,Y]$ 中也是如此.因此 $f_N(Y,X)=C(X,Y)f_N(X,Y)$，从而 $C(X,Y)C(Y,X)=1$，进而 $C(X,Y)=\pm1$.但 $-1$ 不可能，因为这将导致 $f_N(X,X)=0$，于是 $(X-Y)\mid f_N(X,Y)$，与不可约性矛盾.\footnote{译注：原文这一排除反对称情形的论证未单独处理 $N=1$；此时首一极小多项式为 $Y-X$，并不对称.此处保留原论证及其陈述范围.}



再证明 $f_N(X,Y)\in\mathbb Q[X,Y]$.写成 $f_N(X,Y)=\sum c_{i,j}X^iY^j$，其中 $c_{i,j}\in\mathbb C$，且 $c_{i,|\mathcal L|}=\delta_{i,0}$（极小多项式关于 $Y$ 首一）.在方程 $f_N(j,j_N)=0$ 中代入 $j$ 的 $q$ 展开；其系数属于 $\mathbb Z$，故得到整系数线性方程组，并且由极小多项式的唯一性，它确定各 $c_{m,n}$.这个方程组在 $\mathbb C$ 中有解，故在 $\mathbb Q$ 中也有解.所以所有 $c_{m,n}$ 均属于 $\mathbb Q$，从而 $f_N(X,Y)\in\mathbb Q[X,Y]$，如所欲证.



继续之前，需要下面的引理.

\begin{lemma*}

考虑环同态
\begin{equation*}
\begin{tikzcd}[column sep=large,row sep=0.2em]
\mathbb C[X,Y] \arrow[r] & \mathbb C[[q]][q^{-1}][Y]\\
f(X,Y) \arrow[r,mapsto] & f(j,Y),
\end{tikzcd}
\end{equation*}
其中 $j$ 被送到其 $q$ 展开 $j(z)=q^{-1}+\sum_{n\in\mathbb N}c_nq^n$.若 $\mathcal I$ 是 $\mathbb C$ 的对乘法稳定的加法子群，则
\begin{align*}
f(j,Y)\in\mathcal I[[q]][q^{-1}][Y]\quad\Longrightarrow\quad f(X,Y)\in\mathcal I[X,Y].
\end{align*}
\end{lemma*}

\begin{proof}

写成 $f(X,Y)=\sum_i a_i(X)Y^i$，再写 $a_i(X)=\sum_k a_{ik}X^k$.$f(j,Y)$ 中 $q^{-k}Y^i$ 的系数属于 $a_{ik}+\mathbb Z[a_{i,k'}]_{k'>k}$.假设 $f(X,Y)\notin\mathcal I[X,Y]$，并取 $a_{ik}\notin\mathcal I$，使 $k$ 最大.则 $f(j,Y)$ 中 $q^{-k}Y^i$ 的系数属于 $a_{ik}+\mathcal I$，故不属于 $\mathcal I$，矛盾.

\end{proof}

现在，由 $j$ 的 $q$ 展开具有整系数，可知对每个 $\gamma\in\mathcal L$，$j_N(\gamma z)$ 的 $q_N$ 展开系数属于 $\mathbb Z[e^{2i\pi/N}]$.为此，使用引理~\ref{sec:modular-1-3-2}：$\Gamma(1)\begin{psmallmatrix}N&0\\0&1\end{psmallmatrix}\gamma$ 含有唯一的矩阵 $\begin{psmallmatrix}a&b\\0&d\end{psmallmatrix}$，使 $ad=N$ 且 $b<d$.于是 $j_N(\gamma z)=j(\begin{psmallmatrix}a&b\\0&d\end{psmallmatrix}z)$，再利用 $j$ 的 $q$ 展开，就得到所述 $q_N$ 展开.所以 $f_N(j,Y)\in\mathbb Z[e^{2i\pi/N}][[q]][q^{-1}][Y]$，由引理可得 $f_N(X,Y)\in\mathbb Z[e^{2i\pi/N}][X,Y]$.结合 $f_N(X,Y)\in\mathbb Q[X,Y]$，即得
\begin{align*}
f_N(X,Y)\in\mathbb Z[X,Y].
\end{align*}
最后设 $N=p$.回忆第~\ref{sec:modular-1-3-2} 分目中的集合 $\mathcal L(\Gamma(1)\begin{psmallmatrix}p&0\\0&1\end{psmallmatrix}\Gamma(1))$，可见
\begin{align*}
f_N(j,Y)=(Y-j(pz))\prod_{m=0}^{p-1}\left(Y-j\left(\frac{z+m}{p}\right)\right).
\end{align*}
已经看到，$j((z+m)/p)$ 的 $q_p$ 展开系数属于 $\mathbb Z[e^{2i\pi/p}]$.事实上，这些展开模 $\mathbb Z[e^{2i\pi/p}]$ 的理想 $\mathfrak p=(1-e^{2i\pi/p})$ 都同余，因而特别地都同余于 $j_{p^{-1}}(z)=j(z/p)$ 的展开.换言之，在环 $\mathbb Z[e^{2i\pi/p}][[q_p]][q_p^{-1}][Y]$ 中有
\begin{align*}
f_N(j,Y)\equiv(Y-j_p)(Y-j_{p^{-1}})^p\pmod{\mathfrak p}.
\end{align*}
由于两边都属于 $\mathbb Z[[q_p]][q_p^{-1}][Y]$，在后一个环中模 $p$ 也有同样的同余式.在环 $\mathbb Z[[q_p]][q_p^{-1}]$ 中，模 $(p)$ 还有 $j_p\equiv j^p$ 和 $(j_{p^{-1}})^p\equiv j$.于是，在 $\mathbb Z[[q_p]][q_p^{-1}][Y]$ 中有 $f_N(j,Y)\equiv(Y-j^p)(Y^p-j)$.由于这两个元素都属于 $\mathbb Z[[q]][q^{-1}][Y]$，便有
\begin{align*}
f_N(j,Y)-(Y-j^p)(Y^p-j)\in p\mathbb Z[[q]][q^{-1}][Y].
\end{align*}
现在需由此推出
\begin{align*}
f_N(X,Y)-(Y-X^p)(Y^p-X)\in p\mathbb Z[X,Y].
\end{align*}
这由上述引理取 $\mathcal I=p\mathbb Z$ 即得.

\end{proof}

这提示了一个很好的 $\mathbb Q$ 上模型：

\begin{definition*}

将 $\mathbb Q$ 上的模曲线 $X_0(N)_{\mathbb Q}$ 定义为有理函数域为
\begin{align*}
\mathbb Q(j,j_N):=\mathbb Q(j)[j_N]/(f_N(j,j_N))
\end{align*}
的完备代数曲线.

\end{definition*}

根据此定义，将 $X_0(N)_{\mathbb Q}$ 从 $\mathbb Q$ 扩张标量到 $\mathbb C$ 所得的复代数曲线，其对应Riemann曲面就是 $X_0(N)$.也可看出，投影 $X_0(N)\longrightarrow X(1)$ 来自簇的态射 $X_0(N)_{\mathbb Q}\longrightarrow X(1)_{\mathbb Q}$.



\subsection{Galois群在 \texorpdfstring{$\operatorname{Pic}^0(X_0(N))[\ell^m]$}{Pic⁰(X₀(N))[ℓᵐ]} 上的作用}\label{sec:modular-2-3-4}

设 $X_0(N)_{\overline{\mathbb Q}}$ 为 $X_0(N)_{\mathbb Q}$ 扩张标量到 $\overline{\mathbb Q}$ 所得的曲线.除子的概念与Riemann曲面情形完全相同：除子是 $X_0(N)_{\overline{\mathbb Q}}$ 的闭点的和 $\sum_P a_P[P]$.类似地，也有主除子的概念（由于局部环是离散赋值环，可以对有理函数 $f$ 定义 $\operatorname{ord}_P(f)$），继而得到 $\operatorname{Pic}$ 和 $\operatorname{Pic}^0$.



对域 $k$ 上每条光滑完备曲线 $C$，可以赋予其\newterm{modular-jacobian-variety}{Jacobi簇}{Jacobian variety}，这是一个Abel簇，其 $\overline k$ 值点与 $\operatorname{Pic}^0(C_{\overline k})$ 相认同.这一一般事实相当深刻，我们只需它的下列推论：

\begin{enumerate}[label=\roman*)]

\item 有 $\operatorname{Pic}^0(X_0(N)_{\overline{\mathbb Q}})\subset\operatorname{Pic}^0(X_0(N))$；

\item 对每个整数 $M$，有 $\operatorname{Pic}^0(X_0(N)_{\overline{\mathbb Q}})[M]=\operatorname{Pic}^0(X_0(N))[M]$.

\end{enumerate}

换言之，$\operatorname{Pic}^0$ 中先验为复数的挠点都``定义''在 $\overline{\mathbb Q}$ 上.我们现在获得的是 $\operatorname{Gal}(\overline{\mathbb Q}/\mathbb Q)$ 在这些挠点上的作用.事实上，$\operatorname{Gal}(\overline{\mathbb Q}/\mathbb Q)$ 作用于 $X_0(N)_{\overline{\mathbb Q}}$ 的闭点，因而也作用于除子；它显然保持主除子，故作用于 $\operatorname{Pic}^0(X_0(N)_{\overline{\mathbb Q}})$.该作用与群运算相容，因此保持 $M$ 挠点.



接下来需要验证，此Galois作用与Hecke环的作用相容.



\subsection{Hecke对应的有理模型}\label{sec:modular-2-3-5}

回忆，$T(p)$ 的作用就是双陪集
\begin{align*}
\Gamma_0(N)\begin{pmatrix}1&0\\0&p\end{pmatrix}\Gamma_0(N)
\end{align*}
的作用.从几何上看，我们已知它在 $\operatorname{Pic}^0(X_0(N))$ 上的作用由对应
\begin{equation*}
\begin{tikzcd}
X_0(N) & X_0(N,p) \arrow[l,"\pi^p"'] \arrow[r,"\pi"] & X_0(N)
\end{tikzcd}
\end{equation*}
诱导，其中
\begin{align*}
X_0(N,p)=X\left(\Gamma_0(N)\cap\begin{pmatrix}1&0\\0&p\end{pmatrix}^{-1}\Gamma_0(N)\begin{pmatrix}1&0\\0&p\end{pmatrix}\right),
\end{align*}
$\pi$ 是典范投影，而 $\pi^p$ 是 $\begin{psmallmatrix}1&0\\0&p\end{psmallmatrix}$ 的作用与典范投影的复合；这个矩阵诱导同构
\begin{equation*}
\begin{tikzcd}[column sep=small]
{\scriptstyle X\left(\Gamma_0(N)\cap\begin{psmallmatrix}1&0\\0&p\end{psmallmatrix}^{-1}\Gamma_0(N)\begin{psmallmatrix}1&0\\0&p\end{psmallmatrix}\right)}
\arrow[r,"\sim"] &
{\scriptstyle X\left(\begin{psmallmatrix}1&0\\0&p\end{psmallmatrix}\Gamma_0(N)\begin{psmallmatrix}1&0\\0&p\end{psmallmatrix}^{-1}\cap\Gamma_0(N)\right)}.
\end{tikzcd}
\end{equation*}
我们知道，$X_0(N)$ 的函数域是 $\mathbb C(j,j_N)$，因而
\begin{align*}
X\left(\begin{pmatrix}1&0\\0&p\end{pmatrix}^{-1}\Gamma_0(N)\begin{pmatrix}1&0\\0&p\end{pmatrix}\right)
\end{align*}
的函数域为 $\mathbb C(j_{p^{-1}},j_{Np^{-1}})$（注意 $j_{p^{-1}}=j|_0\begin{psmallmatrix}1&0\\0&p\end{psmallmatrix}$ 对 $\begin{psmallmatrix}1&0\\0&p\end{psmallmatrix}^{-1}\Gamma_0(N)\begin{psmallmatrix}1&0\\0&p\end{psmallmatrix}$ 是模函数）.再取域的合成，便得 $\mathcal M_{X_0(N,p)}=\mathbb C(j,j_N,j_{p^{-1}},j_{Np^{-1}})$.于是，用函数域表示，上述对应为
\begin{equation*}
\begin{tikzcd}[column sep=small]
\mathbb C(j,j_N) \arrow[r,"\iota^p"] & \mathbb C(j,j_N,j_{p^{-1}},j_{Np^{-1}}) & \mathbb C(j,j_N) \arrow[l,"\iota"']
\end{tikzcd}
\end{equation*}
其中 $\iota$ 是由 $j\mapsto j$、$j_N\mapsto j_N$ 给出的典范包含映射，而 $\iota^p$ 由 $j\mapsto j_{p^{-1}}$、$j_N\mapsto j_{Np^{-1}}$ 给出.现在显然，同样的映射也诱导图
\begin{equation*}
\begin{tikzcd}[column sep=small]
\mathbb Q(j,j_N) \arrow[r,"\iota^p"] & \mathbb Q(j,j_N,j_{p^{-1}},j_{Np^{-1}}) & \mathbb Q(j,j_N) \arrow[l,"\iota"']
\end{tikzcd}
\end{equation*}
它定义对应 $T(p)$ 的一个 $\mathbb Q$ 上模型，即代数对应
\begin{equation*}
\begin{tikzcd}
X_0(N)_{\mathbb Q} & X_0(N,p)_{\mathbb Q} \arrow[l,"\pi^p"'] \arrow[r,"\pi"] & X_0(N)_{\mathbb Q}.
\end{tikzcd}
\end{equation*}
扩张标量并转到Riemann曲面后，就还原为上述对应.因此，$T(p)$ 在 $\operatorname{Pic}^0(X_0(N)_{\overline{\mathbb Q}})$ 上的作用由Riemann曲面情形下的同一公式 $\pi_*\circ\pi^{p,*}$ 诱导；现在它显然与 $\operatorname{Gal}(\overline{\mathbb Q}/\mathbb Q)$ 的作用交换.

